\justifying
\NSFCBodyText

\subsubsection{与本项目相关的研究工作积累}

申请人颜丰为新疆大学副教授、硕士研究生导师，具备机器学习、表示学习、问答及软件系统研发积累。本课题所需的知识证据组织、检索结果分析与可复现实验，与申请人已有工作存在方法联系。前期研究主要支撑表征、检索及实验实现能力，教育情境标注和智能体风险监测仍是本项目拟开展的研究。

申请人主持的“面向政企数字化应用的智能知识库与问答系统研发”横向课题，提供了知识库建设任务的工程衔接。该课题依合同构建知识库，其客户场景与本项目的教育风险评价对象不同；项目合同中的资料与成果不能未经授权直接作为开放评测数据。教育任务与伦理风险标签将在本课题中独立设计和核验。

\subsubsection{已取得的研究工作成绩}

申请人以第一作者、通讯作者身份完成的Knowledge-aware image understanding with multi-level visual representation enhancement for visual question answering发表于\textit{Machine Learning}，正式卷期为2024年113卷6期3789—3805页\cite{yan2024knowledge}。该工作涉及知识与多层表示的结合，可为本课题组织证据和分析问答行为提供方法经验，不能作为教育智能体可信评价已经有效的实验依据。

申请人作为通讯作者的Feature Fusion Mamba Hashing via Decoupling for Cross-Modal Retrieval发表于\textit{IEEE Signal Processing Letters}，2026年33卷609—613页\cite{tang2026hashing}。申请人在上述成果中负责研究构思、方法设计、实验指导、结果分析及论文撰写与修改，相关贡献已由共用申请人事实材料确认。表示、检索和对照实验方面的积累有助于承担本课题的任务构建及证据分析，但不会把跨模态检索指标移作教育效果指标。

现阶段没有已确认可直接复用的本课题教育轨迹数据、风险评价脚本或消融实验日志。拟建设的EduAgentTrust评测集和监测原型均列为本次研究产出，不作为既有成果展示；因此本申请不提供未经实验支持的性能曲线或前期效果数字。

\subsubsection{可行性分析}

在方法层面，现有通用轨迹与工具风险研究提供了可参考的任务构造、风险标注和基线评价方式；申请人的问答和知识库经验有助于把公开资料组织为可复现任务。在实现层面，已有算力与Linux、Python、深度学习及版本管理工具链能够支持模型推理、日志记录和成对实验。项目以已有模型为对象，不以训练新基础模型为前提。

在组织层面，申请人承担总体方案、技术实现与证据核查；参与本课题的合作人与具体分工为\textbf{【待补 F-EDU-01：中心合作人、团队名单及任务分工】}。教学规则设计和人工评分应由具备相应课程理解能力的人员参与，不能只由算法评价器自行生成并验证标签。如教学规则复核人员无法落实，相关成果限于合成协议验证，不宣称已建立教育风险测量有效性。合作落实前先完成公开资料和沙箱方法准备，以降低场景依赖对总体进度的影响。

\subsubsection{研究风险与应对措施}

\textbf{风险标签存在歧义。}当预实验中分歧集中在帮助程度或教学目标解释时，先补充正反例与任务规则，区分明确风险、支持不足和规范不确定，不通过扩大自动评分规模掩盖分歧。对仍有争议的类别报告敏感性分析，必要时收缩为授权明确的风险任务，保留可复核案例作为最低交付。

\textbf{轨迹评价提升缺少稳定性。}当增益只出现于显式风险词、某一模型或少数模板家族时，检查信息量、长度和任务构造捷径，使用等信息对照与家族留出分析定位原因。若教育约束没有额外作用，按反证结果报告，不通过更换测试集寻求有利结论；将可由规则解决的风险与需要人工判断的风险分别说明。

\textbf{干预降低了正常教学支持。}当误阻断、重规划次数或人工复核负担上升而风险改善有限时，减少自动干预覆盖范围，仅保留确定性权限保护，把教学争议交由人工。评价中保留所有拒绝与失败任务，避免因排除困难样本而高估效果。

\textbf{场景授权或资源进度滞后。}在真实课程数据未获许可时使用公开材料与虚拟角色完成沙箱实验，不接触真实学生身份和成绩。模型接口或成本发生变化时，优先采用可本地复现的模型，保留版本和费用记录；减少双智能体与公平性扩展，确保主要情境对照、评价方法和案例报告在一年内完成。
